\subsection{Reinforcement Learning and Decision Systems}

Reinforcement learning studies agents whose actions affect future observations and rewards.

\begin{definition} \label{def:sl-markov-decision-process} \hlt{Markov Decision Process}\\
A Markov decision process consists of states $\mathcal{S}$, actions $\mathcal{A}$, transition law $P(s'\mid s,a)$, reward function $r(s,a)$, and discount factor $\gamma\in[0,1)$.  A policy $\pi(a\mid s)$ chooses actions from states.
\end{definition}

\begin{definition} \label{def:sl-value-functions} \hlt{Value Functions}\\
The value of a policy is
\[
V^{\pi}(s)=
\mathbb{E}_{\pi}\left[\sum_{t=0}^{\infty}\gamma^t r(S_t,A_t)\,\middle|\,S_0=s\right].
\]
The action-value function $Q^{\pi}(s,a)$ conditions additionally on the first action.
\end{definition}

\begin{definition} \label{def:sl-bellman-equation} \hlt{Bellman Equation}\\
For a fixed policy,
\[
V^{\pi}(s)=
\sum_a\pi(a\mid s)\left(r(s,a)+\gamma\sum_{s'}P(s'\mid s,a)V^{\pi}(s')\right).
\]
Optimality equations replace the policy average by a maximum over actions.
\end{definition}

\begin{remark} \label{rem:sl-rl-learning-systems} \hlt{RL as a Learning System}\\
Reinforcement learning combines prediction, causal feedback, optimization, exploration, and evaluation under distribution shift.  Its validation must account for sequential dependence and policy-induced data.
\end{remark}
